\section{Claim Tiers, Benoit Replication, and Manifesto Provenance}
\label{app:v8-claim-tiers-manifesto}
\label{app:v8-benoit-replication}

This appendix records the evidence tiers and reproduction details behind
Section~\ref{sec:v8-manifesto-results}. No new experiment is introduced here:
the purpose is to make explicit which claims are root-observed, which are
diagnostic, and which would require a completed node-level audit.

\subsection{Claim Tiers}

The v8 manuscript separates the evidence channels because they support
different claims.

\begin{table}[t]
\centering
\small
\begin{tabular}{@{}p{0.22\linewidth}p{0.28\linewidth}p{0.36\linewidth}@{}}
\toprule
Tier & Observation unit & Claim supported \\
\midrule
Root-observed evaluation & Full documents with external targets &
The C-Tree root agrees with the task target on the evaluated corpus. \\
Node-level local-law certificate & Sampled leaves, states, and merge nodes
with logged propensities &
The realized tree satisfies C1/C2/C3 up to estimated oracle distortion. \\
Local supervision & Teacher, rubric, or expert labels on smaller tree units &
The learned \(g\) receives dense task-aligned pressure before or between
oracle audits. \\
Theorem stack & Mathematical assumptions on \(g\), \(T\), and \(f^\ast\) &
Conditional preservation for oracle-compatible readouts. \\
Application narrative & Domain setup and measurement target &
How a concrete task instantiates the state, oracle, and audit units. \\
\bottomrule
\end{tabular}
\caption{\textbf{Claim tiers used in v8.} The tiers can reinforce one another,
but they should not be collapsed.}
\label{tab:v8-claim-tiers}
\end{table}

The Manifesto/Benoit benchmark in the main text is root-observed: each held-out
document has an external full-document expert target, and the tree root is
scored against that target. The node-level local-law certificate is stronger:
it requires sampling realized tree units and evaluating C1/C2/C3 losses through
\(f^\ast\) with logged propensities. The two tiers answer different questions:
the root tier asks whether the final score agrees with experts, while the node
tier asks whether the score is carried by locally faithful policy states.

\subsection{Corpus, Targets, and Expert Benchmarks}

The Manifesto application uses party election programs from the Manifesto
Project matched to expert-survey means from \citet{BenoitEtAl2025}. The
headline benchmark contains \(235\) party election platforms from parliamentary
democracies in election years where Benoit--Laver and CHES expert surveys
provide party-position estimates. Raw text is taken from the Manifesto Corpus
2025a release \citep{ManifestoProject2025a,ManifestoCorpus2025}; the target
labels are the harmonized expert means released in the Benoit replication
archive.

The six dimensions are economic policy, social policy, immigration, European
integration, environmental policy, and decentralization. Pearson \(r\) is
computed across held-out documents for one dimension, and macro \(r\) is the
unweighted average of the six dimension-level correlations. Benoit's
split-expert row is a measurement-agreement reference rather than a training
target: experts are randomly divided into two groups, each group forms a mean
party-position estimate, and the two mean vectors are correlated. The economic
split-expert reference used in the ladder discussion is \(r=0.880\); the macro
reference across the six dimensions is \(0.873\).

The current generated C-Tree artifacts use \(218\) manifestos after local
availability and preprocessing filters. They are split into \(140\) training
documents, \(30\) reserved documents used for model selection and diagnostics,
and \(48\) held-out test documents for final reporting. This \(218\)-document
artifact split is therefore distinct from the \(235\)-document benchmark
definition. The difference is a provenance fact about the generated run, not a
new measurement target.

\subsection{Pipeline Protocol}

Each C-Tree configuration has three components.

\begin{itemize}[leftmargin=1.5em]
    \item \textbf{Leaf summarizer.} Each contiguous text span is summarized
          under a preservation rubric for the active policy dimension or under
          the universal-summary rubric. The prompt targets Benoit's
          \(300\)--\(400\)-word summary regime but allows longer outputs when
          aggressive compression would drop policy-relevant evidence.
    \item \textbf{Merge summarizer.} Adjacent child states are merged under the
          same preservation target. The merge prompt is allowed to carry both
          child summaries forward when that is the least lossy state.
    \item \textbf{Scorer.} The root state is scored with the Benoit rubric for
          dimension \(d\) and mapped to the seven-point policy scale. Invalid
          responses are recorded as missing rather than coerced.
\end{itemize}

The C-Tree rows in Table~\ref{tab:v8-benoit-headline} use an open-weight
\textsc{Gemma-4-31B-IT-NVFP4} model served locally. The exact Benoit scoring
contexts are used where the run is meant to compare directly against the
released Benoit prompts. For the universal-summary run, root states are first
generated for all manifestos and then scored dimension by dimension; this
two-pass schedule keeps the scoring rubric fixed across consecutive calls and
avoids interleaving six long rubric prefixes.

Character-leaf and token-leaf runs are different axes. The main
Manifesto/Benoit table uses the reported character-leaf comparison. The
alternating \(f/g\) ladder reports token leaf sizes because it is generated by
the prompt-optimization artifacts. The two should not be compared as if they
were one continuous leaf-size sweep.

\subsection{Controls and Robustness Diagnostics}

The full per-cell grid behind the main table includes direct and local
controls. A flat baseline scores raw manifesto text directly after truncation;
across truncation sizes it gives macro \(r\) in the \(0.786\)--\(0.827\) band.
This is the long-context measurement control: it measures how much of the
expert-survey signal is visible to a direct read before the document is
distributed over tree leaves. A concat baseline chunks the document and scores
the concatenated leaf summaries without recursive merges; it gives macro
\(r\) around \(0.827\)--\(0.837\) at the strongest reported settings. The
concat row shows that the first local compression step preserves much of the
benchmark signal before the tree merge operator is introduced.

Decoding sensitivity was checked by low-temperature and medium-temperature
runs, with and without three-sample aggregation. At the \(8\mathrm{K}\)
character setting, the per-dimension C-Tree macro lies in
\([0.832,0.841]\) across these variants, while the universal-summary C-Tree
lies in \([0.800,0.808]\). The rank ordering of the main configurations is
therefore not driven by one decoding temperature.

The older replication materials also include an exact-model Gemma-3 control on
Benoit's anonymized summaries. Running the scorer on those summaries with the
Gemma-3-27B BF16 model used in Benoit's Table 6 produces macro \(r\) in the
\(0.723\)--\(0.731\) band across rubric and prompt-format variants. This
control is used only to document model-version sensitivity; it is not the
reported v8 C-Tree result.

Robustness slices are consistent with the main interpretation. Among
manifestos longer than \(100\mathrm{K}\) characters, the per-dimension C-Tree
reports Pearson \(r=0.89\) on economic policy, \(0.85\) on social policy,
\(0.87\) on immigration, \(0.88\) on European integration, and \(0.83\) on
environment. Decentralization remains weaker at \(0.47\), matching the main
table rather than creating a new long-document failure mode. Election-era
slices show the same pattern: economic, social, environmental, immigration,
and European integration remain high across available periods, while
decentralization stays lower.

Benoit et al. also discuss Manifesto Project quasi-sentence counts transformed
into coded logit scores. We treat those as a construct-validity check, not as
the target benchmark here, because they measure a different object from
expert-survey party-position means. In the generated checks, tree, flat,
concat, and universal variants produce substantially lower macro correlations
against the coded-score target than against expert means. That gap is evidence
that the expert-survey target and coded-count target should not be treated as
interchangeable validation labels.

\subsection{Decentralization as the Stress Dimension}

Decentralization is the dimension where both Benoit and the C-Tree artifacts
show the strongest warning signal. In Benoit's expert benchmark it has a lower
split-expert reference than the other five dimensions. In the C-Tree artifacts
it is the dimension most sensitive to leaf size and shared-state compression.
Older fixed-prompt sweeps show decentralization improving as leaves shrink:
\(r=0.489\) at \(64\mathrm{K}\) characters, \(0.543\) at \(8\mathrm{K}\), and
\(0.584\) at \(2\mathrm{K}\). The v8 interpretation is that finer leaves force
the tree to carry institutional qualifications, regional claims, and
central-local tradeoffs into the merge state instead of smoothing them away in
one large leaf summary.

This pattern is why the main text treats decentralization as a local-law
diagnostic rather than as a generic model-quality failure. The root score
identifies the weak dimension, but the local diagnostic gap identifies where
new node labels would be most valuable.

\subsection{Compute and Call Accounting}

The cost comparison is reported as call accounting, not as a universal price
claim.

\begin{itemize}[leftmargin=1.5em]
    \item \textbf{Benoit proprietary ensemble.} Per
          manifesto--dimension, the released ensemble uses
          \(3\) LLMs \(\times\) \(2\) prompt settings \(\times\)
          \(3\) summaries, or \(18\) calls. Across \(235\) manifestos and
          \(6\) dimensions this is \(25{,}380\) LLM calls routed through
          frontier APIs.
    \item \textbf{Universal-summary C-Tree.} Per manifesto, the run uses
          \(2\lceil n_{\mathrm{chars}}/c\rceil-1\) summarize/merge calls plus
          \(6\) scoring calls. At \(c=8\mathrm{K}\) and a typical
          \(40\)--\(48\)K-character manifesto, this is about \(15\)--\(17\)
          calls on one local open-weight model.
    \item \textbf{Per-dimension C-Tree.} Six dimension-specific trees are
          built per manifesto. At \(c=8\mathrm{K}\), this is roughly
          \(60\)--\(70\) calls per manifesto on the same local model.
\end{itemize}

Dollar costs depend on hardware ownership, utilization, and API pricing, so
the manuscript does not claim a precise head-to-head price ratio. The
reproducible point is the call structure: one local open-weight model with a
tree state map versus a proprietary ensemble with multiple frontier models and
multiple prompt settings.

\subsection{Optimization-Stage Provenance}
\label{app:v8-rile-fg-ladder}

The \(f/g\) stage labels in Section~\ref{sec:v8-manifesto-results} come from
alternating prompt-program search on the scorer \(f\) and state map \(g\).
Model weights are not updated; the optimized object is the prompt program. A
stage label \(f^a g^b\) means that the scorer-side prompt has been updated
\(a\) times and the summarizer-side prompt has been updated \(b\) times. It is
a stage label, not multiplication.

Two initialization lanes are tracked. The Benoit-init lane starts from
archived Benoit summaries, so its anchor is already a scorer-optimized
\(f^1g^0\) state. The raw-init lane starts from the local Gemma baseline
summary and an untuned scorer, so the first step isolates scorer adaptation.
When a later stage reused the same evaluation as the preceding stage, the
table reports only the first stage at which the distinct number appeared.

The Benoit-init lane reaches an external \(r\) plateau of
\([0.879,0.886]\) from \(1024\) tokens onward, straddling the economic
split-expert reference of \(0.880\). The raw-init lane contains the strongest
single economic cell, \(r=0.909\), at \(256\) tokens after the first
summarizer update. These ladder numbers are not compared directly to the
per-dimension headline row because the leaf axis and optimization budget differ.
The full generated grid is kept as
\path|assets/benoit/tables/manifesto_fg_ladder.tex| for provenance, but v8
does not render it as a second results table.

\subsection{Asset Provenance}

\begin{table}[t]
\centering
\small
\begin{tabular}{@{}p{0.32\linewidth}p{0.52\linewidth}@{}}
\toprule
Artifact & Use in v8 \\
\midrule
\path|assets/benoit/tables/benoit_comparison_pearson.tex| &
Main six-dimension Manifesto/Benoit comparison,
Table~\ref{tab:v8-benoit-headline}. \\
\path|assets/benoit/figures/manifesto_singledim_per_dim_live.pdf| &
Single-dimension held-out trajectories,
Figure~\ref{fig:v8-benoit-singledim-per-dim}. \\
\path|assets/benoit/figures/manifesto_fg_combined_ladder_f1g0_f1g1.pdf| &
Universal-summary held-out trajectories,
Figure~\ref{fig:v8-benoit-combined-ladder}. \\
\path|assets/benoit/figures/manifesto_fg_combined_audit_gap.pdf| &
Universal-summary teacher-trace diagnostic gaps,
Figure~\ref{fig:v8-benoit-combined-audit-gap}. \\
\path|assets/benoit/figures/manifesto_singledim_per_dim_live_audit_gap.pdf| &
Single-dimension teacher-trace diagnostic gaps,
Figure~\ref{fig:v8-benoit-single-audit-gap}. \\
\path|assets/benoit/tables/manifesto_fg_ladder.tex| &
Generated economic \(f/g\) stage grid retained for provenance; not rendered as
a second v8 results table. \\
\path|assets/benoit/tables/length_buckets.md| and
\path|assets/benoit/tables/era_breakdown.md| &
Robustness diagnostics summarized in this appendix. \\
\bottomrule
\end{tabular}
\caption{\textbf{Manifesto artifact provenance.} The current v8 manuscript
treats these generated files as fixed artifacts; it documents them rather than
regenerating them.}
\label{tab:v8-benoit-asset-provenance}
\end{table}

\subsection{Oracle-Grounded Local-Law Path}

The present Manifesto result is meaningful as a root-observed evaluation. The
stronger local-law path is to sample realized leaves, state re-entry calls,
and merge nodes, then evaluate their C1/C2/C3 losses against \(f^\ast\). In
the Manifesto domain, expert quasi-sentence codings or high-quality teacher
judgments are natural candidates for such local oracle calls, provided they
measure the same policy dimension as the root target.

The audit design must declare the span aggregation rule, the active law
channels, and the logged inclusion propensities before reporting the
finite-population estimate. Quasi-sentence labels are therefore presented here
as the natural route to node-level oracle evidence, not as already-used node
audit evidence. Teacher traces and \(\widehat f\) provide dense local
supervision for \(g\); sampled \(f^\ast\) node calls estimate
\(b_v=\widehat\ell_v-\ell_v^\ast\) and correct the local-law channel without
changing the root/local population objective.
